% TODO make hw stylesheet
\documentclass[10pt]{article}
\usepackage[letterpaper,top=1in,left=1in,right=1in]{geometry}
\usepackage[dvipsnames,svgnames]{xcolor}
\usepackage[shortlabels]{enumitem}
\usepackage[framemethod=TikZ]{mdframed}
\usepackage{amsmath,amssymb,amsthm}
\usepackage[colorlinks]{hyperref}
\usepackage{multicol}
\usepackage{tabularx}
\usepackage{bbm}

\hypersetup{
    colorlinks=true,
    linkcolor=blue,
    filecolor=magenta,
    urlcolor=magenta,
    pdftitle={MAT 324 HW}
}

\usepackage{mathtools}
\usepackage{thmtools}
\usepackage[textsize=footnotesize]{todonotes}
\usepackage{titling}
\usepackage{cancel}

\reversemarginpar
\mdfdefinestyle{mdbluebox}{roundcorner=10pt,innerbottommargin=9pt,
linecolor=blue,backgroundcolor=TealBlue!5,}
\declaretheoremstyle[headfont=\sffamily\bfseries\color{MidnightBlue},
mdframed={style=mdbluebox},]{thmbluebox}

\mdfdefinestyle{mdbloobox}{frametitlefont=\bfseries,innerbottommargin=8pt,
nobreak=true,backgroundcolor=TealBlue!5,linecolor=blue,}
\declaretheoremstyle[headfont=\bfseries\color{MidnightBlue},
mdframed={style=mdbloobox},headpunct={\\[3pt]},postheadspace=0pt,]{thmbloobox}

\mdfdefinestyle{mdredbox}{frametitlefont=\bfseries,innerbottommargin=8pt,
nobreak=true,backgroundcolor=Salmon!5,linecolor=RawSienna,}
\declaretheoremstyle[headfont=\bfseries\color{RawSienna},
mdframed={style=mdredbox},headpunct={\\[3pt]},postheadspace=0pt,]{thmredbox}

\mdfdefinestyle{mdgreenbox}{linecolor=ForestGreen,backgroundcolor=ForestGreen!5,
linewidth=2pt,rightline=false,leftline=true,topline=false,bottomline=false,}
\declaretheoremstyle[headfont=\bfseries\sffamily\color{ForestGreen!70!black},
mdframed={style=mdgreenbox},headpunct={ --- },]{thmgreenbox}

\mdfdefinestyle{mdorangebox}{linecolor=RedOrange,backgroundcolor=RedOrange!5,
linewidth=2pt,rightline=false,leftline=true,topline=false,bottomline=false,}
\declaretheoremstyle[headfont=\bfseries\sffamily\color{RedOrange!70!black},
mdframed={style=mdorangebox},headpunct={ --- },]{thmorangebox}

\mdfdefinestyle{mdblackbox}{linecolor=black,backgroundcolor=RedViolet!5!gray!5,
linewidth=3pt,rightline=false,leftline=true,topline=false,bottomline=false,}
\declaretheoremstyle[mdframed={style=mdblackbox}]{thmblackbox}

\declaretheoremstyle[spaceabove=6pt,spacebelow=6pt]{boldhead}
\declaretheoremstyle[spaceabove=6pt,spacebelow=6pt,headfont=\normalfont\itshape]{italichead}

\declaretheorem[style=boldhead,name=Problem]{problem}
\declaretheorem[style=boldhead,name=Problem,numbered=no]{problem*}
\declaretheorem[style=boldhead,name=Setup]{setup} % for config geo
\declaretheorem[style=boldhead,name=Example,sibling=problem]{example}
\declaretheorem[style=boldhead,name=$\bigstar$ Problem,sibling=problem]{niceprob}
\declaretheorem[style=boldhead,name=Theorem,sibling=problem]{theorem}
\declaretheorem[style=boldhead,name=Corollary,sibling=theorem]{corollary}
\declaretheorem[style=boldhead,name=Proposition,sibling=theorem]{prop}
\declaretheorem[style=boldhead,name=Lemma,numberwithin=problem]{lemma}
\declaretheorem[style=boldhead,name=Theorem,numbered=no]{theorem*}
\declaretheorem[style=boldhead,name=Lemma,numbered=no]{lemma*}
\declaretheorem[style=boldhead,name=Claim,numberwithin=problem]{claim}
\declaretheorem[style=boldhead,name=Claim,numbered=no]{claim*}
\declaretheorem[style=boldhead,name=Remark,numberwithin=problem]{remark}
\declaretheorem[style=boldhead,name=Remark,numbered=no]{remark*}
\declaretheorem[style=boldhead,name=Definition,sibling=theorem]{definition}
\declaretheorem[style=boldhead,name=Definition,numbered=no]{definition*}

\providecommand{\ol}{\overline}
\providecommand{\ceil}[1]{\left\lceil #1 \right\rceil}
\providecommand{\floor}[1]{\left\lfloor #1 \right\rfloor}
\newcommand{\dd}{\,\mathrm{d}}
\providecommand{\ul}{\underline}
\providecommand{\eps}{\varepsilon}
\providecommand{\half}{\frac{1}{2}}
\providecommand{\CC}{\mathbb C}
\providecommand{\FF}{\mathbb F}
\providecommand{\II}{\mathbb I}
\providecommand{\NN}{\mathbb N}
\providecommand{\QQ}{\mathbb Q}
\providecommand{\RR}{\mathbb R}
\providecommand{\ZZ}{\mathbb Z}
\providecommand{\ind}{\mathbbm 1}
\providecommand{\dg}{^\circ}
\providecommand{\ii}{\item}
\providecommand{\setdiff}{\smallsetminus}
\providecommand{\alert}[1]{{\sffamily\textbf{\textcolor{blue}{#1}}}}
\providecommand{\inv}{^{-1}}
\providecommand{\mb}[1]{\mathbf{#1}}

\newenvironment{soln}{\begin{proof}[Solution]}{\end{proof}}

\providecommand{\pow}{\operatorname{Pow}}
\providecommand{\vocab}[1]{{\textbf{\textcolor{ForestGreen}{#1}}}}
\providecommand{\lcm}{\operatorname{lcm}}
\providecommand{\abs}[1]{\left\lvert #1\right\rvert}
\providecommand{\innerprod}[1]{\langle\!\langle #1\rangle\!\rangle}
\providecommand{\norm}[1]{\left\| #1\right\|}
\providecommand{\smnorm}[1]{\| #1\|}
\providecommand{\gr}{\operatorname{gr}}

\usepackage{mathrsfs}

% place title at top right
\setlength{\droptitle}{-8ex}
\pretitle{\begin{flushright}\Large\bfseries}
\posttitle{\par\end{flushright}}
\preauthor{\begin{flushright}}
\postauthor{\end{flushright}}
\predate{\begin{flushright}}
\postdate{\end{flushright}}

\title{MAT 324 HW10}
\author{\large Aditya Pahuja}
\date{\large Due 11/20}
\begin{document}
\maketitle

\begin{problem}
    [Axler 8A \#17]
    Let $\lambda$ denote Lebesgue measure on $[1,\infty)$.
    \begin{enumerate}[(a)]
	\ii Prove that if $f\colon[1,\infty)\to[0,\infty)$
	is Borel measurable, then
	\begin{equation*}
	    \left(\int_{[1,\infty)}f\dd\lambda\right)^2
	    \le \int_{[1,\infty)} x^2f(x)^2\dd\lambda(x).
	\end{equation*}
	\ii Describe the set of Borel measurable functions
	$f\colon[1,\infty)\to[0,\infty)$ such that
	the inequality in (a) is an equality.
    \end{enumerate}
\end{problem}

\begin{soln}
    If the left-hand side is $\infty$, then so is the right-hand side because\dots;
    \todo{deal with the LHS being infinite}
    and if the right-hand side is infinite then the inequality is trivial.
    Otherwise, $f$ and $x^2 f(x)^2$ are both in
    the $L^2$-space $L^2([1,\infty); \RR_{\ge0})$
    of nonnegative square-integrable functions,
    so by Cauchy-Schwarz on the functions $xf(x)$ and $1/x$,
    \begin{equation*}
	\left(\int_{[1,\infty)} xf(x)\cdot\frac 1x\dd\lambda(x)\right)^2
	\le \left(\int_{[1,\infty)} x^2 f(x)^2\dd\lambda(x)\right)
	{\underbrace{\left(\int_{[1,\infty)}
	\frac 1{x^2}\dd\lambda(x)\right)}_{=1/1 - 1/\infty = 1}}.\qedhere
    \end{equation*}
\end{soln}

\begin{problem}
    Let $(X,\mathcal S,\mu)$ be a measure space
    where there exist disjoint $A,B\in\mathcal S$ of positive, finite measure.
    Show that the norm $\norm{\bullet}_p$ on $L^p(\mu)$ is not induced
    by an inner product on $L^p(\mu)$ for any $p\in[1,\infty]-\{2\}$.
\end{problem}

\begin{soln}
    If $\norm{\bullet}_p$ is induced by an inner product,
    then the parallelogram law holds for the $\mathcal S$-measurable functions
    $\ind_A,\ind_B\colon X\to\RR$, i.e.
    \begin{equation*}
	\norm{\ind_A + \ind_B}_p^2 + \norm{\ind_A - \ind_B}_p^2
	= 2\norm{\ind_A}_p^2 + \norm{\ind_B}_p^2.
    \end{equation*}
    Since $A$ and $B$ are disjoint,
    $\abs{\ind_A + \ind_B} = \abs{\ind_A - \ind_B} = \ind_{A\cup B}$,
    so the sum and difference of the two indicator functions
    have the same $L^p$-norm. Thus, after dividing out a factor of two,
    \begin{equation*}
	(\mu(A) + \mu(B))^{2/p} = \mu(A\cup B)^{2/p}
	= \mu(A)^{2/p} + \mu(B)^{2/p}.\tag{$\ast$}
    \end{equation*}

    \begin{lemma}
        \label{lem:nonadditivity}
        Let $x,y,p\in\RR^+$. The equality $x^p + y^p = (x + y)^p$ holds
	if and only if $p = 1$.
    \end{lemma}
    \begin{proof}
        Let $t = x/y$, so the given equality rewrites as
	$t^p + 1 = (t+1)^p$ upon dividing by $y^p$.
	Then, we want to show that $f_p(t) = (t+1)^p - t^p - 1$
	has no positive zeros. Note that $f_p'(t) = p(t+1)^{p-1} - pt^{p-1}$;
	we analyze the sign of the derivative based on the value of $p$.
	\begin{itemize}
	    \ii If $p-1 > 0$, then $(t+1)^{p-1} > t^{p-1}$ for all $t\ge 0$,
	    meaning $f$ is increasing on $\RR_{\ge0}$.
	    Since $f_p(0) = 0$, $f_p$ is positive for all positive $t$.
	    \ii If $p-1 < 0$, then $(t+1)^{p-1} < t^{p-1}$ for all $t > 0$,
	    meaning $f_p$ is decreasing on $\RR^+$.
	    Since $f_p$ is continuous on $\RR_{\ge0}$, it is
	    decreasing on the closure $\RR_{\ge0}$ of $\RR^+$,
	    so, since $f_p(0) = 0$, $f_p$ is negative for all positive $t$.
	\end{itemize}
	Thus if $p\ne 1$, the equality doesn't hold,
	and of course if $p = 1$ then $f_p(t) = 0$ everywhere.
    \end{proof}

    Applying the lemma with $x = \mu(A)$ and $y = \mu(B)$,
    we get that $(\ast)$ holds if and only if $2/p = 1$, hence
    the parallelogram law doesn't hold if $p \ne 2$,
    so $\norm{\bullet}_p$ is not be induced by an inner product for $p\ne 2$.
\end{soln}

\begin{problem}
    Let $C^{0,2}(\RR)\subseteq L^2(\RR)$ denote
    the subspace of continuous square-integrable functions. Define
    \begin{equation*}
	L_0^2(\RR) = \{f\in L^2(\RR) : f = 0\text{ a.e. on }[0,1]\},\qquad
	C_0^{0,2}(\RR) = L_0^2(\RR)\cap C^{0,2}(\RR).
    \end{equation*}
    \begin{enumerate}[(a)]
	\ii Let $f\in C^{0,2}(\RR)$ be such that $f(x)\ne 0$
	for some $x\in\RR\setminus[0,1]$.
	Show that there exists $g\in C_0^{0,2}(\RR)$ such that
	$\langle\!\langle f,g \rangle\!\rangle_2\ne 0$.
	Conclude that $f\in C^{0,2}(\RR)$ has a projection to $C^{0,2}_0(\RR)$
	if and only if $f(0) = f(1) = 0$.
	\ii Let $f\in L^2(\RR)$. Determine the projection of $f$ to $L_0^2(\RR)$.
    \end{enumerate}
\end{problem}

\begin{soln}
    (a) Suppose $f(x_0)\ne 0$ for $x\in\RR\setminus[0,1]$;
    we can assume without loss of generality that $f(x_0) > 0$.
    Since $f$ is continuous at $x_0\in\RR\setminus[0,1]$,
    there is an open interval $I\subseteq\RR\setminus[0,1]$ containing $x_0$
    on which $f$ is positive.
    Then, let $g\colon\RR\to[0,1]$ be a continuous function
    whose support is contained in $I$ and for which $g(x_0) = 1$.
    This means $g\in C_0^{0,2}(\RR)$, $f(x)g(x) = 0$ for $x\in\RR\setminus I$,
    and since $f(x_0)g(x_0) > 0$, $\int_I fg\dd\lambda > 0$, so
    \begin{equation*}
	\innerprod{f,g}_2 = \int_{\RR} fg\dd\lambda
	= \int_I fg\dd\lambda +
	\underbrace{\int_{\RR\setminus I}fg\dd\lambda}_{=0}> 0.
    \end{equation*}

    Now, we deal with the existence of the projection of $f$ onto $C_0^{0,2}(\RR)$.
    \begin{lemma}
	The distance from $f$ to $C_0^{0,2}(\RR)$ is $\smnorm{f|_{[0,1]}}_2$.
    \end{lemma}
    \begin{proof}
	For $\delta>0$, let $h_\delta\in C_0^{0,2}(\RR)$ be defined as
	\begin{equation*}
	    h_\delta(x) = \begin{cases}
		0 & \text{if }x\in[0,1]\\
		f(x) & \text{if }x\in\RR\setminus(-\delta,1+\delta)\\
		x\cdot \frac{f(-\delta)}{-\delta} & \text{if }x\in(-\delta,0)\\
		(x-1)\cdot \frac{f(1+\delta)}{\delta} & \text{if }x\in(1,1+\delta).
	    \end{cases}
	\end{equation*}
	This is easily checked to be a continuous function
	that agrees with $f$ everywhere except on $(-\delta,1+\delta)$.
	Letting $M$ be the maximum of $\abs f$ over $\RR$
	(which exists because $f^2$ has finite integral),
	$h_\delta$ is at most $2M$ on the intervals
	$(-\delta,0)$ and $(1,1+\delta)$, so
	\begin{equation*}
	    \norm{f-h_\delta}_2^2 =
	    \int_\RR\abs{f-h_\delta}^2\dd\lambda
	    \le \int_{(-\delta,0)\cup(1,1+\delta)} \abs{f-h_\delta}^2\dd\lambda
	    + \int_{[0,1]} \abs{f}^2\dd\lambda
	    \le (2M)^2\delta + \norm{f\cdot\ind_{[0,1]}}_2^2.
	\end{equation*}
	As $\delta$ goes to zero, we get that the distance from $f$ to
	$C_0^{0,2}(\RR)$ is bounded above by $\smnorm{f|_{[0,1]}}_2$.
	Of course, $\smnorm{f\cdot\ind_{[0,1]}}_2$ is also a lower bound
	on the distance because
	$\norm{f-h}_2^2 = \int_\RR\abs{f-h}^2\dd\lambda
	\ge \int_{[0,1]} \abs{f}^2\dd\lambda$
	whenever $h$ is zero on $[0,1]$; this completes the proof.
    \end{proof}

    To finish, $h_f\in C_0^{0,2}(\RR)$ is the projection of $f$ onto $C_0^{0,2}(\RR)$
    if and only if
    \begin{equation*}
	\norm{f-h_f}_2^2 = \smnorm{f\cdot\ind_{[0,1]}}^2
	\iff \int_{\RR\setminus[0,1]} \abs{f-h_f}^2\dd\lambda = 0.
    \end{equation*}
    The latter equality is equivalent to $f-h_f$ being the zero function,
    which is equivalent to $f-h_f$ being zero
    on the closure $\RR\setminus(0,1)$ of $\RR\setminus[0,1]$.
    Since $f-h_f$ and $f$ agree at $0$ and $1$, such an $h_f$
    can only exist if $f(0) = f(1) = 0$,
    and indeed if $f(0) = f(1) = 0$ then $h_f = f\cdot\ind_{[0,1]}$
    is a continuous function that
    minimizes the distance between $f$ and a point of $C_0^{0,2}(\RR)$.
    \\

    (b) The distance from $f\in L^2(\RR)$ to $L_0^2(\RR)$
    can be bounded below by $\smnorm{f\cdot\ind_{[0,1]}}_2$ because
    \[
	\norm{f-h}_2^2 = \int_\RR\abs{f-h}^2\dd\lambda
	\ge \int_{[0,1]} \abs{f}^2\dd\lambda
    \]
    whenever $h$ is zero on $[0,1]$,
    and the lower bound is achieved by taking $h$
    to be $h_f = f\cdot\ind_{[0,1]}$, so this $h_f$ is the point
    minimizing the distance from $f$ to a point of $L_0^2(\RR)$,
    which is by definition the projection.
\end{soln}

\begin{problem}
    Let $\mathbb I = [0,1]$ and $L^2(\mathbb I;\CC)$ be the $L^2$-space
    of (equivalence classes of) $\CC$-valued functions on $\mathbb I$.
    For each $n\in\ZZ$ and $f\in L^2(\mathbb I,\CC)$, define
    \begin{equation}
	\psi_n\colon\mathbb I\to\CC,\;\;\psi_n(x) = e^{2\pi i nx},\qquad
	c_n(f) = \innerprod{f,\psi_n}_2
	= \int_{\mathbb I} f\ol{\psi_n}\dd\lambda\in\CC.
	\label{eq:p4-assumptions}
    \end{equation}
    \begin{enumerate}[(a)]
	\ii Show that the collection $\{\psi_n : n\in\ZZ\}$
	consists of orthonormal elements of $L^2(\mathbb I;\CC)$.
	\ii Let $f\in L^2(\mathbb I;\CC)$. Show that the sum
	\begin{equation*}
	    \sum_{n\in\ZZ} c_n(f)\psi_n\coloneq
	    \lim_{k,m\to\infty}\sum_{n=-k}^m c_n(f)\psi_n
	\end{equation*}
	converges with respect to the $L^2$ norm
	to an $h_f\in L^2(\mathbb I,\CC)$ so that
	$\norm{h_f}_2\le\norm{f}_2$ and
	$\innerprod{f-h_f,\psi_n}_2 = 0$ for all $n\in\ZZ$.
	\ii Suppose $f$ is twice continuously differentiable,
	$f(0) = f(1)$, and $f'(0) = f'(1)$. Show that for all $n\ne 0$,
	\begin{equation*}
	    c_n(f) = \frac 1{2\pi in}c_n(f')
	    = -\frac 1{4\pi^2 n^2}c_n(f'').
	\end{equation*}
	\ii Under the assumptions on $f$ in (c), show that the sum in (b)
	converges uniformly to a continuous function $h_f\colon\mathbb I\to\CC$.
\end{enumerate}
\end{problem}

\begin{soln}
    (a) If $m=n$ then $\psi_m(x)\ol{\psi_n(x)} = e^{2\pi i(m-n)x}$
    is $1$ for all $x\in\II$, so $\innerprod{\psi_m,\psi_m}_2 = 1$
    for all $m\in\ZZ$, and if $m\ne n$ then
    \begin{equation*}
	\innerprod{\psi_m, \psi_n}_2 =
	\int_{\II} e^{2\pi i(m-n)x}\dd\lambda(x)
	= \frac{1}{2\pi i(m-n)}(e^{2\pi i(m-n)} - e^0) = 0.
    \end{equation*}

    (b) Axler's 8.57 implies that
    \[
	\sum_{n\in\ZZ} \abs{c_n(f)}^2 =
	\sum_{n\in\ZZ}\norm{c_n(f)\psi_n}_2^2 \le \norm{f}_2^2,
    \]
    which by Axler's 8.54 implies that $\sum_{n\in\ZZ}c_n(f)\psi_n$ converges
    in the $L^2$ norm to some $h_f$ and
    \begin{equation*}
	\norm{\sum_{n\in\ZZ} c_n(f)\psi_n}_2^2
	= \sum_{n\in\ZZ} \abs{c_n(f)}^2 \le \norm{f}_2^2.
    \end{equation*}
    Furthermore, by Axler's 8.58(b),
    $h_f\in\ol{\operatorname{span}\{\psi_n : n\in\ZZ\}}$
    is equal to $\sum_{n\in\ZZ}\innerprod{h_f,\psi_n}\psi_n$
    \begin{equation*}
	\innerprod{f-h_f,\psi_n}_2 = \innerprod{f,\psi_n}_2 -
	\sum_{m\in\ZZ}\innerprod{c_m(f)\psi_m,\psi_n}_2
	= \innerprod{f,\psi_n}_2 - c_n(f) = 0.
    \end{equation*}
    \todo{technicality with taking inf sum out of inner prod}

    (c) For any continuously differentiable function $g$
    with $g(0) = g(1)$, integration by parts yields
    \begin{equation*}
	c_n(g) = \int_{\II} g\ol{\psi_n}\dd\lambda
	= \frac 1{2\pi i(-n)}(g(1)\ol{\psi_n(1)} - g(0)\ol{\psi_n(0)})
	- \int_{\II}g'\cdot\frac 1{2\pi i(-n)}\ol{\psi_n}\dd\lambda
	= \frac 1{2\pi in}c_n(g').
    \end{equation*}
    Applying this to $g = f$ and $g = f'$ gives us what we want.
    \\

    (d) For each nonzero integer $n$,
    \begin{equation*}
	\abs{c_n(f)} = \frac{1}{4\pi^2n^2}\abs{c_n(f'')}
	\le \frac 1{4\pi^2n^2}\int_{\II}\abs{f''\ol{\psi_n}}\dd\lambda
	= \frac 1{4\pi^2n^2}\norm{f''}_1,
    \end{equation*}
    so, since $\sum_{n=1}^\infty \frac1{4\pi^2n^2} < \infty$,
    \begin{equation*}
	\sum_{n\in\ZZ}\abs{c_n(f)} \le \abs{c_0(f)}
	+ 2\sum_{n=1}^\infty\frac 1{4\pi^2n^2}\norm{f''}_1
	< \infty,
    \end{equation*}
    which implies that the sum in (b) converges uniformly to $h_f$
    by the Weierstrass $M$-test
    since $\abs{c_n(f)\psi_n} = \abs{c_n(f)}$ for all $x$.
    The summands $c_n(f)\psi_n$ are all continuous,
    so this implies $h_f$ is continuous as well.
\end{soln}

\begin{problem}
    Let $\mathbb I = [0,1]$, $L^2(\mathbb I,\CC)$ be
    the $L^2$-space of $\CC$-valued functions on $\mathbb I$,
    $f\in L^2(\mathbb I,\CC)$, and $\psi_n$ and $c_n(f)$ be as in
    (\ref{eq:p4-assumptions}).
    \begin{enumerate}[(a)]
	\ii Suppose $f$ is continuous and $f(0) = f(1)$.
	Show that, given $\eps>0$, there exists an $N\in\NN$
	and $a_n\in\CC$ for $n=-N,\dots,N-1,N$ such that
	\begin{equation*}
	    \abs{f(x) - \sum_{n=-N}^N a_n\psi_n(x)}\le \eps
	    \quad\text{for all }x\in\mathbb I.
	\end{equation*}
	\ii Suppose $f$ is twice continuously differentiable,
	$f(0) = f(1)$, and $f'(0) = f'(1)$. Show that
	\begin{equation*}
	    f = \sum_{n\in\ZZ} c_n(f)\psi_n
	\end{equation*}
	with the sum converging uniformly on $\mathbb I$.
	\ii Show that the sum in (b) converges to $f$
	with respect to the $L^2$ norm for every $f\in L^2(\mathbb I,\CC)$.
\end{enumerate}
\end{problem}

\begin{soln}
    (a) 
    \\

    (b)
    \\

    (c)
\end{soln}










\end{document}
